\documentclass[../cours_2de.tex]{subfiles}
\begin{document}
\setcounter{chapter}{13}
\chapter{Valeur absolue et fonction carré}
\label{sec:ch14}
\begin{objectif}
Connaître deux courbes de référence ; démontrer signes et variations ; résoudre des équations et inéquations avec des distances ou des carrés.
\end{objectif}
\section{La fonction valeur absolue}
$f(x)=|x|$ est définie sur $\R$ par $f(x)=-x$ pour $x\leq0$ et $f(x)=x$ pour $x\geq0$.
Sa courbe est formée de deux demi-droites, et non d'une droite unique.
Elle est positive ou nulle, s'annule seulement en $0$, décroît sur $]-\infty,0]$ et croît sur $[0,+\infty[$.
\begin{demonstration}
Sur chaque demi-droite, les expressions sont affines de pentes respectives $-1$ et $1$. Le chapitre 13 donne les variations. La définition comme distance donne le signe et le minimum $0$ atteint en $0$.
\end{demonstration}
On a $|-x|=|x|$ : deux points opposés sont à la même distance de l'origine.
Plus généralement, une fonction est \emph{paire} si son domaine est symétrique par rapport à $0$ et si $f(-x)=f(x)$ ; sa courbe est alors symétrique par rapport à l'axe des ordonnées.
\begin{demonstration}
Le symétrique de $(x,f(x))$ par rapport à l'axe vertical est $(-x,f(x))$.
Il appartient à la courbe exactement lorsque $f(-x)=f(x)$.
\end{demonstration}
\begin{exemple}
$|x-3|=2$ décrit les deux points à distance $2$ de $3$ : $x=1$ ou $x=5$.
$|x-3|<2$ donne $]1,5[$.
$|x-3|>2$ donne $]-\infty,1[\cup]5,+\infty[$.
Si le second membre est négatif, une égalité à une valeur absolue n'a aucune solution.
\end{exemple}
\section{La fonction carré}
La fonction $g(x)=x^2$ est définie sur $\R$. Elle est paire, positive ou nulle, et s'annule seulement en $0$. Sa courbe est une parabole.
\begin{theoreme}
La fonction carré est strictement décroissante sur $]-\infty,0]$ et strictement croissante sur $[0,+\infty[$. Son minimum vaut $0$, atteint seulement en $0$.
\end{theoreme}
\begin{demonstration}
Si $0\leq a<b$, alors $b^2-a^2=(b-a)(b+a)>0$.
Si $a<b\leq0$, alors $b-a>0$ mais $a+b<0$, donc $b^2-a^2<0$.
Enfin, un carré est non négatif et ne vaut zéro que si sa base vaut zéro.
\end{demonstration}
\begin{center}
\begin{tikzpicture}
\begin{axis}[width=10cm,height=5.8cm,axis lines=middle,xmin=-2.3,xmax=2.3,ymin=-.3,ymax=4.5,xlabel={$x$},ylabel={$y$},legend pos=north west,samples=100]
\addplot[accent,thick,domain=-2.1:2.1]{x^2};\addlegendentry{$x^2$}
\addplot[dashed,domain=-2.1:2.1]{abs(x)};\addlegendentry{$|x|$}
\end{axis}
\end{tikzpicture}
\end{center}
\section{Résoudre avec un carré}
Pour $k>0$, $x^2=k$ donne $x=\pm\sqrt k$.
Comme $x^2$ dépend de la distance de $x$ à $0$ :
\[
x^2<k\Leftrightarrow-\sqrt k<x<\sqrt k,\qquad
x^2\geq k\Leftrightarrow x\leq-\sqrt k\text{ ou }x\geq\sqrt k.
\]
Pour $k=0$, traiter séparément l'égalité et le sens de l'inégalité ; pour $k<0$, utiliser la positivité du carré.
\begin{exemple}
$(x-2)^2\leq9$ équivaut à $-3\leq x-2\leq3$, donc à $x\in[-1,5]$.
$x^2>-1$ est vraie pour tous les réels ; $x^2<0$ n'a aucune solution.
\end{exemple}
\section{Comparer les courbes de $x$ et $x^2$}
\begin{demonstration}
Pour $x\geq0$, la différence $x^2-x=x(x-1)$ est négative sur $]0,1[$, nulle en $0$ et $1$, positive pour $x>1$.
La parabole est donc sous la droite entre $0$ et $1$, et au-dessus après $1$.
\end{demonstration}
Ainsi $0{,}2^2<0{,}2$, mais $3^2>3$. « Le carré agrandit toujours » est faux.
Pour $x<0$, le carré est non négatif tandis que $x$ est négatif : $x^2>x$.
\section{Une parabole déplacée}
\approfondissement Pour $F(x)=(x-\alpha)^2+\beta$, le minimum est $\beta$, atteint en $\alpha$.
La fonction décroît jusqu'à $\alpha$ puis croît ; sa courbe a pour axe $x=\alpha$.
Cela résulte du changement de variable $u=x-\alpha$, qui conserve l'ordre des entrées, suivi de l'ajout de $\beta$.
\begin{exemple}
$x^2-6x+11=(x-3)^2+2$.
On en déduit $F(x)\geq2$, et $F(x)=6$ équivaut à $(x-3)^2=4$, donc à $x=1$ ou $x=5$.
La méthode vient de l'identité remarquable ; une formule générale du second degré n'est pas nécessaire ici.
\end{exemple}
\utilite{Le carré intervient dans les aires, les distances et la dispersion statistique. Comprendre ses variations évite les erreurs lorsqu'on élève une inégalité au carré.}
\begin{retenir}
Sur des nombres négatifs, le carré renverse l'ordre. Une équation carrée peut avoir deux solutions ; une distance se traite en tenant compte des deux côtés du centre.
\end{retenir}
\end{document}
